\section{Temporal records and SI mapping}
\label{manual:m3-temporal}
\paragraph{Constructed input.}
The adapter consumes the admitted recurrent waves and repeated coupled returns of
\mainequation{main:eq:wave-update-trace},
\mainequation{main:eq:coupling-trace} and
\mainequation{main:eq:rcd-additive-realization}.
It retains each complete trace, certificate, lineage, contact prefix, integer
window endpoint, signed phase bias, lag and directional reflection count.
The worked pair has four compared closures at cuts \(2,10,18,26\) and three
reference closures at cuts \(2,12,22\). Over \((0,30]\), excluding the first
closure of each recurrent lineage, its increments are \((3,2)\).
The comparison ratio is \(3/2\), its reverse is \(2/3\), and the two
additive directional durations are \(5\) and \(7\) in their declared native
count domain. Biases \((1/4,-1/4)\) give lag \(3/2\) and reverse lag \(-3/2\).

\paragraph{Observation domain.}
The mapping binds the population, preparation, directed couplings, complete
committed prefixes, pending returns and window convention. Directional RCD, bias, lag and drift are recomputed from those records.
Structural RD belongs to each retained support. Duration participation requires
a compatible window map; a reclosure count alone does not supply it. Displacement traversals, hinges, dwells, emitted ticks and
returned clicks retain their separate counts. A window with positive reclosure
counts can still lack the repeated cross-returns needed for RCD.

\paragraph{Declared observation and units.}
An observation law first maps reference history \(\Gamma_R\), window
\(\omega\), population \(\mathcal P\), coupling history \(E\) and preparation
to corrected Cs-referenced cycle equivalents:
\begin{equation}
 \widehat C_{\mathrm{Cs}}
   =g_\theta(\Gamma_R,\omega,\mathcal P,E,\mathrm{preparation}),
 \qquad
 \widehat T=\frac{\widehat C_{\mathrm{Cs}}}{\Delta\nu_{\mathrm{Cs}}}.
 \label{eq:manual-m3-cycle-bridge}
\end{equation}
The conversion uses the exact defining frequency; the first map requires
its own observation and calibration evidence.
Let \(a_c,a_r>0\) be stipulated cycle-equivalent gains per native return.
Let \(f_{\mathrm{Cs}}=9\,192\,631\,770\,\mathrm{Hz}\) denote the defining
reference, without a clock realization. For the fixed window define
\[
 T=\frac{a_rN_r}{f_{\mathrm{Cs}}},\qquad
 \bar\nu_c=\frac{a_cN_c}{T},\qquad
 t_1=t_0+T.
\]
The mapping binds both gains, the time coordinate and epoch, the complete
input records, the defining reference and the executable mapping specification.
At \(a_c=a_r=1\), \(T=2/f_{\mathrm{Cs}}\) seconds and
\(\bar\nu_c=3f_{\mathrm{Cs}}/2\).
The result is a finite-window average.
The reciprocal \(1/\bar\nu_c\) is a derived period, distinct from \(T\)
and from an independent validation. Physical realization and uncertainties from
drift, endpoint phase, cycle slips and temporal dependence remain unspecified.
The gains are stipulated without fitting hydrogen.

\paragraph{Identifiability and separated evaluation.}
The family \(g_\lambda(Z)=1+\lambda(55-Z)/54\), \(\lambda>-1\), has
\(g_\lambda(55)=1\) and \(g_\lambda(1)=1+\lambda\).
The \(\lambda=0,1\) countermodels preserve the native records and Cs anchor
while producing different reports. The relative gain therefore remains free.
The known hydrogen value remains retrospectively exposed. Evaluation validates
the frozen prediction against the original model inputs before comparison,
then binds separate prediction and target hashes. Changing only the
target changes its diagnostic residual and leaves prediction bytes unchanged.
An atomic point prediction and its complete uncertainty remain unresolved.

\paragraph{Frozen transfer test.}
The retained direct-coupling experiment calibrates on its first half and
validates on the second. Each half contains \(32\) reference reclosures.
With \(k_i=32/N_{i,\mathrm{cal}}\) frozen, the predicted validation count is
\(\widehat N_{0,\mathrm{val}}=k_iN_{i,\mathrm{val}}\).
\begin{table}[ht]
\centering
\caption{Frozen gains and subsequent reference-count transfer.}
\label{tab:manual-m3-anchor-transfer}
\begin{tabular}{lrrrr}
\toprule
Field & \(N_{i,\mathrm{cal}}\) & \(N_{i,\mathrm{val}}\) &
 \(\widehat N_{0,\mathrm{val}}\) & \((\widehat N_{0,\mathrm{val}}-32)/32\)\\
\midrule
A & 58 & 48 & \(768/29\) & \(-5/29\)\\
B & 151 & 102 & \(3264/151\) & \(-49/151\)\\
\bottomrule
\end{tabular}
\end{table}
These residuals reject this constant-transfer specialization. They are
count-transfer discrepancies, not measured Cs-frequency errors or standard
uncertainties. Refitting on each validation window would remove the test.
Identical synchronized reference histories do not provide independent
evidence. Overlapping windows and shared calibration sources retain their
dependence.

\paragraph{Population and reference controls.}
Four retained executions use two recurring \(3{:}4{:}6\) background supports
and the same frozen counting windows.
The baseline compares \(3{:}3{:}6\) with \(1{:}2{:}6\); a second construction compares
\(3{:}2{:}6\) with \(3{:}3{:}6\). Its A counts change from \(99\) to \(72\), while
its reference changes from \(31\) to \(30\). Two separate interventions
activate a prepared \(3{:}4{:}6\) local support after the common committed prefix.
Coupling it only to A gives A counts \(45\) then \(42\), with reference
counts \(32\) in both windows. Coupling it only to the reference gives A
counts \(57\) then \(47\), with reference counts \(23\) in both windows.
Complete histories retain directed reflections and pending returns; short
count projections do not acquire RCD admission. These are finite-model
controls, without measured cycle calibration.

For a physical test, corrected observed cycle increments must be supplied
separately and compared using a justified tolerance fixed before evaluation.
The oscillator, counter, interrogation, realization corrections, endpoint
resolution and missed-cycle model belong to that observation record.
No isotope retention parameter absorbs the residual.
Each seconds or hertz output keeps its own eight-component budget; missing
physical calibration and discrepancy remain unknown.

\paragraph{Consumer and causal control.}
The separately bound Matter I v2 replay preserves its declared-model status:
its native atomic admission remains pending. The baseline and capacity
intervention share committed history through cut 39. The intervention at 40
causes retention loss at 41, changing observed increments from \((6,5)\) to
\((2,5)\). Applying the same stipulated gain one in both cases leaves
\(T=5/f_{\mathrm{Cs}}\) fixed and changes the average from
\(6f_{\mathrm{Cs}}/5\) to \(2f_{\mathrm{Cs}}/5\).
A hostile mutation disconnects the loss response from the next state while
retaining the intervention inputs. The effect disappears and the original
replay validator rejects that trace.

\paragraph{Reproduction and source precision.}
\texttt{manual/scripts/aod\_m3\_temporal.py --check} reconstructs the native
adapter and mapping report; \texttt{--consumer-root} selects the
separately bound Matter I replay.
\texttt{aod\_m3\_anchor.py --check-data} replays the original anchor and
the four population controls. The tests reject changed scopes, aliases,
uncoupled records, invalid windows, scheduler substitutions, nonpositive gains,
and evaluation leakage. The consumer preserves historical v1/v2 fixture bytes
and identifies its two active Main amendments separately.
The BIPM comparison retains raw hashes, tests nearest rounding for nonzero
entries and uses conditional point-zero covariance transcriptions.
Its source documents do not specify zero formatting or underflow, so zero
precision remains unresolved.

\aodliteraturenote{The BIPM defines the second through
\(\Delta\nu_{\mathrm{Cs}}=9\,192\,631\,770\,\mathrm{Hz}\)~\cite{bipm-second}.
In a fountain clock, microwave interrogation and detected atomic-state
populations guide oscillator tuning; an electronic counter counts microwave
cycles~\cite{nist-fountain}. These measured quantities require an observation
map before comparison with native reclosures.}
